\uuid{TZGK}
\titre{Convergence}
\auteur{bodin ; hanani ; mzari}
\organisation{exo7}
\licence{CC BY-NC-SA 4.0}
\datecreate{2020-01-20}
\niveau{L1}
\module{Analyse}
\chapitre{Suite}
\sousChapitre{Autre}
\difficulte{2}
\theme{Suites, convergence}
\exo7qcmid{lille-1.suites.016}

\contenu{
\begin{question}
\qclassification{121.00}
\qtags{motcle=convergence}
Soit $\displaystyle u_n=\sin\left(\frac{2n\pi}{3}\right)$ et $\displaystyle v_n=\sin\left(\frac{3}{2n\pi}\right)$. Quelles sont les bonnes réponses ?
\begin{answers}
    \good{La suite $(u_n)$ diverge et la suite $(v_n)$ converge.}
    \bad{Les suites $(u_n)$ et $(v_n)$ sont divergentes.}
    \good{La suite $(u_n)$ n'a pas de limite et $\displaystyle \lim _{n\to +\infty}v_n=0$.}
    \bad{$\displaystyle \lim _{n\to +\infty}u_n=0$ et $\displaystyle \lim _{n\to +\infty}v_n=0$}
\end{answers}
\begin{explanations}
Par continuité de la fonction sinus, on a :
$$\displaystyle \lim _{n\to +\infty}\frac{3}{2n\pi}=0\mbox{ donc }\displaystyle \lim _{n\to +\infty}\sin\left(\frac{2n\pi}{3}\right)=\sin\left(\lim _{n\to +\infty}\frac{3}{2n\pi}\right)=\sin 0=0.$$
Ainsi, $(v_n)$ converge et sa limite est $0$. Par ailleurs, $$u_{3n}=\sin (2n\pi)=0\mbox{ et }u_{3n+1}=\sin \left(\frac{2\pi}{3}\right)=\frac{\sqrt{3}}{2}.$$
Donc, d'après le théorème des suites extraites, $(u_n)$ diverge ; elle n'a pas de limite.
\end{explanations}
\end{question}
}
